% Part: first-order-logic
% Chapter: axiomatic-deduction
% Section: rules-and-proofs

\documentclass[../../../include/open-logic-section]{subfiles}

\begin{document}

\iftag{FOL}
      {\olfileid{fol}{axd}{rul}}
      {\olfileid{pl}{axd}{rul}}

\olsection{قواعد و \usetoken{P}{derivation}}

\begin{explain}
  دستگاه‌های !!{derivation} اصل‌موضوعی شاید ساده‌ترین دستگاه‌های
  !!{derivation} برای منطق باشند. !!^a{derivation} صرفاً دنباله‌ای از
  !!{formula}s است. برای آنکه دنباله‌ای !!a{derivation} به شمار آید، هر
  !!{formula} در آن باید یا نمونه‌ای از یک اصل موضوع باشد، یا به‌وسیلهٔ
  یک قاعدهٔ استنتاج از یک یا چند !!{formula}s پیشین در دنباله به دست آید.
  !!^a{derivation} با !!{derive}s، آخرین !!{formula} خود را نتیجه می‌دهد.
\end{explain}

\begin{defn}[!!^{derivability}]
اگر $\Gamma$ مجموعه‌ای از !!{formula}s زبان $\Lang L$ باشد، آنگاه
\emph{!!{derivation} از}~$\Gamma$ دنبالهٔ متناهی
$!A_1$، \dots،~$!A_n$ از !!{formula}s است که برای هر $i \le n$ یکی از
شرایط زیر برقرار باشد:
\begin{enumerate}
\item $!A_i \in \Gamma$؛ یا
\item $!A_i$ یک اصل موضوع است؛ یا
\item $!A_i$ به‌وسیلهٔ یک قاعدهٔ استنتاج از یک $!A_j$ (و $!A_k$) با
  $j < i$ (و $k < i$) به دست می‌آید.
\end{enumerate}
\end{defn}

اینکه چه چیزی !!{derivation} درست به شمار می‌آید، به قواعد استنتاجی که
مجاز می‌دانیم (و البته به آنچه اصل موضوع می‌گیریم) بستگی دارد. قاعدهٔ
استنتاج گزاره‌ای شرطی است که می‌گوید گام~$!A_i$ در !!a{derivation} تحت
چه شرایطی گام استنتاجیِ درستی است.

\begin{defn}[قاعدهٔ استنتاج]
\emph{قاعدهٔ استنتاج} شرطی کافی برای آنچه گام استنتاجی درست در
!!a{derivation} از~$\Gamma$ به شمار می‌آید به دست می‌دهد.
\end{defn}

برای نمونه، چون هر دنبالهٔ تک‌عضوی $!A$ با $!A \in \Gamma$ بدیهتاً
!!a{derivation} به شمار می‌آید، عبارت زیر می‌تواند قاعدهٔ استنتاج بسیار
ساده‌ای باشد:
\begin{quote}
  اگر $!A \in \Gamma$، آنگاه $!A$ همواره گام استنتاجی درستی در هر
  !!{derivation} از~$\Gamma$ است.
\end{quote}
همچنین، اگر $!A$ یکی از اصول موضوع باشد، آنگاه $!A$ به‌تنهایی
!!a{derivation} است؛ پس این نیز قاعدهٔ استنتاج است:
\begin{quote}
  اگر $!A$ یک اصل موضوع باشد، آنگاه $!A$ گام استنتاجی درستی است.
\end{quote}
هنگامی که قاعدهٔ استنتاج به !!{formula}s پیش از گام مورد نظر استناد کند،
موضوع جالب‌تر می‌شود. قاعدهٔ زیر \emph{وضع مقدم:} نام دارد
\begin{quote}
  اگر $!B \lif !A$ و $!B$ بالاتر در !!{derivation} آمده باشند، آنگاه
  $!A$ گام استنتاجی درستی است.
\end{quote}
اگر این تنها قاعدهٔ استنتاج باشد، تعریف بالا از !!{derivation} به این
معناست: $!A_1$، \dots،~$!A_n$~!!a{derivation} است هرگاه و تنها هرگاه
برای هر $i \le n$ یکی از شرایط زیر برقرار باشد:
\begin{enumerate}
\item $!A_i \in \Gamma$؛ یا
\item $!A_i$ یک اصل موضوع است؛ یا
\item برای یک $j < i$، $!A_j$ همان $!B \lif !A_i$ باشد و برای یک
  $k < i$، $!A_k$ همان~$!B$ باشد.
\end{enumerate}
بند آخر می‌گوید $!A_i$ با وضع مقدم از~$!A_j$ ($!B \lif !A_i$) و
$!A_k$ ($!B$) به دست می‌آید. اگر بتوانیم از $1$ تا~$n$ پیش برویم و هر
بار !!a{formula}~$!A_i$ بیابیم که یا در~$\Gamma$ است، یا اصل موضوع است،
یا قاعدهٔ استنتاجی به ما می‌گوید گام استنتاجی درستی است، آنگاه کل دنباله
!!{derivation} درست به شمار می‌آید.

\begin{defn}[!!^{derivability}]
!!{formula}~$!A$~\emph{!!{derivable}} از $\Gamma$ است، و
$\Gamma \Proves !A$ نوشته می‌شود، اگر !!a{derivation} از $\Gamma$ وجود
داشته باشد که به $!A$ پایان یابد.
\end{defn}

\begin{defn}[قضیه‌ها]
!!^a{formula}~$!A$ یک \emph{قضیه} است اگر !!a{derivation} از مجموعهٔ
تهی برای~$!A$ وجود داشته باشد. $\Proves !A$ می‌نویسیم اگر $!A$ قضیه
باشد و اگر نباشد، $\Proves/ !A$.
\end{defn}


\end{document}
